% !TITLE 蜀相
% !AUTHOR 杜甫
% !DYNASTY 唐
% !ORDER 26

\begin{poem}
丞相祠堂何处寻，锦官城外柏森森\textsuperscript{1}。
映阶碧草自春色\textsuperscript{2}，隔叶黄鹂空好音\textsuperscript{3}。
三顾频烦天下计\textsuperscript{4}，两朝开济老臣心\textsuperscript{5}。
出师未捷身先死\textsuperscript{6}，长使英雄泪满襟。
\end{poem}

\begin{annotations}
\notetext{1}{柏森森：古柏茂密的样子。武侯祠前种着柏树，四季常青，象征忠贞。}
\notetext{2}{自春色：草自己绿着。不管有没有人欣赏，春光自顾自地美好——反衬人的心情凄凉。}
\notetext{3}{空好音：鸟叫得再好听也是白费。和“自春色”一样，用自然的美好反衬内心的失落。}
\notetext{4}{三顾频烦天下计：刘备三顾茅庐，多次向诸葛亮请教统一天下的计策。频烦，屡次。}
\notetext{5}{两朝开济老臣心：诸葛亮辅佐刘备（先主）和刘禅（后主）两朝，开创基业、匡济危局，一片老臣忠心。开，开创；济，匡济。}
\notetext{6}{出师未捷身先死：出兵北伐还没有胜利，自己在五丈原先去世了。指234年诸葛亮最后一次北伐，病逝军中。}
\end{annotations}

\begin{appreciation}
《蜀相》是杜甫最著名的咏史诗之一。写的是诸葛亮，其实写的也是他自己。

丞相祠堂何处寻，锦官城外柏森森——注意这个"寻"字。不是路过看见了，是专程来找。锦官城就是成都，武侯祠今天还在。祠堂前古柏森森——柏树四季常青，种在祠堂前，取其忠贞不渝的意思。

映阶碧草自春色，隔叶黄鹂空好音——草自己绿着，鸟自己叫着，跟祠堂里的人、来凭吊的人都没关系。"自"和"空"是全诗的机关：春天的美好是客观的，心情的凄凉是主观的，两下一撞，美好反而成了讽刺。这就是以乐写哀。还记得《采薇》里"昔我往矣，杨柳依依"吗？春天越好，走的人越难受，是同一个道理。

三顾频烦天下计，两朝开济老臣心——十四个字，写完诸葛亮的一生：隆中对策，赤壁东风，白帝托孤，六出祁山。开是开创，济是匡济。老臣心三个字最重：他不是为功名，是认定了的事，就做到底。这让我想起《离骚》里的屈原——"亦余心之所善兮，虽九死其犹未悔"，认定了，九死也不后悔。诸葛亮也是这种人。

出师未捷身先死，长使英雄泪满襟——建兴十二年，诸葛亮最后一次北伐，病死在五丈原军中，五十四岁。"长使"两个字值得咂摸：为什么长久地让英雄落泪？因为英雄代代都有，代代都有人想做诸葛亮那样的事，代代都有人等不到结局。杜甫哭诸葛亮，一半是敬，一半是哭自己。他一辈子想济世安民，结果安史之乱一打八年，大半生在逃难。站在武侯祠前，他看见的不是古人，是一个想做事而不得的自己。

这首诗写在成都草堂刚建成不久。天下不安，个人的安稳就总带着亏欠——这份心思，从《采薇》的战士到《离骚》的屈原，再到杜甫，一脉相承。

你以后做事，多半也会遇到"出师未捷"的时候。到时候别急着丧气，把该做的做完，结果好坏，再说。
\end{appreciation}
